\documentclass{article}
\usepackage[T1]{fontenc}
\usepackage{lmodern}
\usepackage{graphicx}
\usepackage{amsmath,amssymb}
\usepackage[a4paper,margin=2cm]{geometry}
\usepackage[absolute,overlay]{textpos}
\setlength{\TPHorizModule}{1mm}
\setlength{\TPVertModule}{1mm}
\usepackage[indonesian]{babel}
\usepackage{hyperref}
\setlength{\emergencystretch}{2em}
% unit: Halaman 86

\begin{document}

% === Halaman 86 (llm) ===

\begin{verbatim}
void salsa20_block(uint32_t out[16], uint32_t const in[16])
{
    int i;
    uint32_t x[16];

    for (i = 0; i < 16; ++i)
        x[i] = in[i];
    // 10 loops x 2 rounds/loop = 20 rounds
    for (i = 0; i < ROUNDS; i += 2) {
        // Column round
        QR(x[0], x[4], x[8], x[12]); // column 1
        QR(x[5], x[9], x[13], x[1]); // column 2
        QR(x[10], x[14], x[2], x[6]); // column 3
        QR(x[15], x[3], x[7], x[11]); // column 4
        // Row round
        QR(x[0], x[1], x[2], x[3]); // row 1
        QR(x[5], x[6], x[7], x[4]); // row 2
        QR(x[10], x[11], x[8], x[9]); // row 3
        QR(x[15], x[12], x[13], x[14]); // row 4
    }
    for (i = 0; i < 16; ++i)
        out[i] = x[i] + in[i];
}
\end{verbatim}

\begin{flushright}
Sumber: Wikipedia
\end{flushright}

\end{document}
